\documentclass[12pt]{article}

\usepackage{url,verbatim, amsmath, amssymb, amsfonts, dsfont, color}
\textwidth 15true cm
\textheight 9true in
\oddsidemargin 0.30true in
\evensidemargin 0.30true in
\topmargin -0.3true in
\headsep 0.4true in

\def\sgn{{\rm sign}}
\def\drightarrow{{\stackrel{d}{\rightarrow}}}
\def\dotsim{\stackrel{.}{\sim}}
\def\goesd{\stackrel{d}{\rightarrow}}
\def\goesp{\stackrel{p}{\rightarrow}}
\def\goeslaw{\stackrel{{\cal L}}{\longrightarrow}}
%
\def\vp{{\varphi}}
\def\E{{\rm E}}
\def\Pvalue{{$p$-value}}
\def\Pvalues{{$p$-values}}
\def\pr{{\rm pr}}
\def\Pr{{\rm pr}}
\def\logit{{\rm logit}}
\def\T{{ \mathrm{\scriptscriptstyle T} }}
\def\diag{{\rm diag}}  
\def\half{{\textstyle{1\over 2}}}
\def\ith{{$i{\rm th}$ }}
\def\var{{\rm var}}
\def\cl{{{\rm c}\ell}}
\def\checkx{{\checked\hspace{-6pt}\setminus}}
\def\grad{{\nabla}}
\def\bs#1{{\boldsymbol{#1}}}

\newcommand{\R}{\mathbb{R}}
\newcommand{\N}{\mathbb N}
\newcommand{\Z}{\mathbb Z}
\newcommand{\Q}{\mathbb Q}
\newcommand{\1}{\mathds 1}
\newcommand{\Var}{\text{Var}}
\newcommand{\Cov}{\text{Cov}}
\newcommand{\sd}{\text{sd}}
\newcommand{\se}{\text{se}}
\renewcommand{\L}{\mathcal {L}}
\renewcommand{\lambda}{\uplambda}
\newcommand{\iidsim}{\stackrel{iid}{\sim}}
\newcommand{\otherwise}{\text{otherwise}}
\newcommand{\indep}{\perp \!\!\! \perp}
\renewcommand{\epsilon}{\varepsilon}
\newcommand{\st}{\text{ s.t. }}
\newcommand{\ds}{\displaystyle}
\newcommand{\ind}{\stackrel{d}{\to}}
\newcommand{\inp}{\stackrel{p}{\to}}
\newcommand{\inas}{\stackrel{a.s.}{\to}}
\DeclareMathOperator*{\argmax}{arg\,max}
\DeclareMathOperator*{\argmin}{arg\,min}

\sloppy %permits line-breaking of urls

\def\blue{\color{blue}}
\def\red{\color{red}}
\def\sgn{{\rm sign}}
\def\drightarrow{{\stackrel{d}{\rightarrow}}}
\begin{document}

\noindent{\bf Questions Week 7} {\hfill STA 2212S 2026}\\

%{\bf Due January 14} \hfill{\it MS, Exercise 5.2}

\bigskip

\begin{enumerate}

\item In Davison's analysis of publication bias,\footnote{Davison, A.C. {\it Statistical Models}, Chapter 5.5} he considers a model in which a single study with $n$ individuals leads to estimate $\hat\mu \dotsim N(\mu, \sigma^2/n)$; that this study is published, if some measure of randomness in publication, denoted $Z$, is positive, 
	and that $\hat\mu$ and $Z$ are related according to the model  $$\hat\mu = \mu + \sigma n^{-1/2} U_1, \quad Z = \gamma_0 + \gamma_1n^{1/2} + U_2, \quad \text{cor}(U_1,U_2)=\rho \ge 0.$$ 
	\begin{enumerate}
	\item Show that $\pr(R=1) = \pr(Z > 0) = \Phi(\gamma_0+\gamma_1n^{1/2})$ and 
	$$\displaystyle{\pr(R=1\mid\hat\mu) = \pr(Z > 0 \mid \hat\mu) = \Phi\left\{\frac{\gamma_0 + \gamma_1n^{1/2} + \rho n^{1/2}(\hat\mu-\mu)/\sigma)}{(1-\rho^2)^{1/2}}\right\}}.$$  
	\item  Show that the estimate of $\mu$ is biased if $\rho>0$: \begin{align*}E(\hat\mu \mid R=1) &= \mu+\rho\sigma n^{-1/2}\zeta(\gamma_0+\gamma_1 n^{1/2}), %\mu+\rho\sigma\gamma_1\zeta'(\gamma_0)+\rho\sigma\zeta(\gamma_0)n^{-1/2}
	\end{align*} where $\zeta(x) = \phi(x)/\Phi(x)$, and $\phi(\cdot), \Phi(\cdot)$ are the standard normal density and distribution function, respectively
	\item For a collection of $k$ such studies, each producing estimate $\hat\mu_j \dotsim N(\mu, \sigma^2/n_j)$, derive the log-likelihood function given in Davison's eq.(5.34): $$ \ell(\mu,\rho, \sigma^2) = \sum_{j=1}^k \left\{\frac{1}{2}\log\sigma^2 + \frac{n_j}{2\sigma^2}(\hat\mu-\mu)^2 - \log\Phi(a_j)+\log\Phi(b_j) \right\}, $$
	where $a_j = \gamma_0+\gamma^1n_j^{1/2}$ and $ b_j = \{a_j+\rho n_j^{1/2}(\hat\mu_j-\mu)/\sigma\}(1-\rho^2)^{-1/2}$.
	\item Verify that if $\rho=0$ that the maximum likelihood estimate of $\mu$ is $$\frac{\sum_{j=1}^k n_j\hat\mu_j}{\sum_{j=1}^k n_j}.$$
\end{enumerate}
\item{\it Davison, Exercise 5.8, Applying the EM algorithm to censored data}
\begin{enumerate}
	\item Suppose $U$ is an underlying failure time, and that $Y=\min(U, c)$ and $D = I(U\le c)$ are observed, where $c$ is a fixed censoring time. The complete data log-likelihood function is $f(u;\theta)$. Show that $$f(u\mid y,d;\theta) = \begin{cases} \delta(u-y),&d=1, \\ \frac{f(u;\theta)}{1-F(c;\theta)}, & u>c, d=0, \end{cases}$$ where $\delta(u-y)$ is  the generalized derivative of the Heaviside function $H(u) = 1\{u\ge 0\}$, and satisfies $\int \delta(y-u)g(u)du = g(y)$ for any function $g(\cdot)$.\footnote{$\delta(\cdot)$ is called the Dirac delta function.}
	\item If $f(u;\theta) = \theta e^{-\theta u}$, show that $\E_{\theta'}(U\mid Y=y, D=d) = dy +(1-d)(c+1/\theta')$, and deduce that the estimate of $\theta$ at the $j$th iterate of the EM algorithm is $$\theta^{(j)} = \frac{n}{\sum_{j=1}^n\{d_jy_j + (1-d_j)(c+1/\theta^{ (j-1)})\}}.$$
	\item The basis of the EM algorithm is summarized by Eq.(5.36) in Davison: $$\log f(y,u;\theta) = \log f(y;\theta) + \log f(u\mid y;\theta);$$ here $U$ are the unobserved (latent) variables and the LHS is called the complete-data likelihood function.   Derive the equality 
    \begin{align*}
         -\frac{\partial^2\log f(y;\theta)}{\partial\theta\partial\theta^T} = \frac{\partial^2\ell(\theta)}{\partial\theta\partial\theta^T} &= \E_\theta\left\{-\frac{\partial^2\log f(y,U;\theta)}{\partial\theta\partial\theta^T}\mid Y=y \right\} \\ 
	    & ~~~~~~~~~~-\E_\theta\left\{-\frac{\partial^2\log f(U\mid y;\theta)}{\partial\theta\partial\theta^T}\mid Y=y\right\},
	\end{align*}
	i.e.~ $J(\theta) = I_c(\theta;y) - I_m(\theta;y)$, interpreted as meaning that the observed information equals the complete-data information minus the missing information.
	\item Show that the missing information for the model in (b) is $\sum_{j=1}^k(1-d_j)/\theta^2$. 
\end{enumerate}
\textcolor{red}{
(a) Suppose that $d = 1$. Then 
\begin{equation*}
	f(u \mid y, d = 1 ; \theta) = \delta(u-y)
\end{equation*}
since $u = y$ by definition. When $d = 0$ and $u > c$, 
\begin{align*}
	f(u \mid y, d = 0 ; \theta) &= \frac{P(Y - y, d = 0 \mid u)f(u)}{P(Y = y, \delta = 0)} \\
	&= \frac{\1_{U >c} f(u)}{1 - F(c)}
\end{align*}
hence 
\begin{equation*}
	f(u \mid y, d ; \theta) = \begin{cases}
		\delta(u - y) & d = 1\\
		\frac{f(u ; \theta)}{1 - F(c ; \theta)} & u > c , d = 0
	\end{cases}
\end{equation*}
}

\textcolor{red}{
(b) 
\begin{align*}
	E_{\theta'}[U \mid Y = y, D = d] &= d E_{\theta'}[U \mid Y = y, D = 1] + (1 - d) E_{\theta '} [U \mid Y = y, \delta = 0]\\
	&= d y + (1 - d) E_{\theta'} [U \mid U > c]\\
	&= d y + (1 - d) \left(c + \frac 1{\theta'}\right)
\end{align*}
}
\end{enumerate}
\end{document}

